\documentclass{article}
\usepackage{booktabs}
\usepackage{longtable}
\usepackage{array}
\usepackage{geometry}
\usepackage{xcolor}
\usepackage{textcomp}
\usepackage{amsmath}
\geometry{margin=1in}

\begin{document}

\title{ANOVA Assumptions Assessment for Box-Cox Transformed Habitat Features}
\author{Taxa Assemblage Groups Analysis}
\date{\today}
\maketitle

\section{Main Results Table}

\begin{longtable}{p{2.8cm}ccccccc}
\caption{Assessment of ANOVA Assumptions for Box-Cox Transformed Environmental Features} \\
\toprule
\textbf{Feature} & \textbf{$\lambda$} & \textbf{Shapiro-W} & \textbf{p-value} & \textbf{Normality} & \textbf{Levene's} & \textbf{p-value} & \textbf{Homog.} \\
\midrule
\endfirsthead

\multicolumn{8}{c}{\tablename\ \thetable\ -- \textit{Continued from previous page}} \\
\toprule
\textbf{Feature} & \textbf{$\lambda$} & \textbf{Shapiro-W} & \textbf{p-value} & \textbf{Normality} & \textbf{Levene's} & \textbf{p-value} & \textbf{Homog.} \\
\midrule
\endhead

\midrule
\multicolumn{8}{r}{\textit{Continued on next page}} \\
\endfoot

\bottomrule
\endlastfoot

Temperature (\textdegree{}C) & 1.902 & 0.812 & 0.000 & \textcolor{red}{FAIL} & 7.978 & 0.000 & \textcolor{red}{FAIL} \\
Depth (m) & 0.239 & 0.986 & 0.642 & \textcolor{green}{PASS} & 0.970 & 0.412 & \textcolor{green}{PASS} \\
DO Bottom (mg/L) & 2.611 & 0.907 & 0.000 & \textcolor{red}{FAIL} & 0.656 & 0.582 & \textcolor{green}{PASS} \\
LOI (\%) & 0.060 & 0.993 & 0.958 & \textcolor{green}{PASS} & 2.657 & 0.056 & \textcolor{green}{PASS} \\
MPS (Phi) & 3.304 & 0.920 & 0.000 & \textcolor{red}{FAIL} & 0.611 & 0.610 & \textcolor{green}{PASS} \\

\end{longtable}

\section{Summary Statistics}

\begin{table}[h]
\centering
\caption{Summary of Assumption Test Results: Box-Cox vs Original}
\begin{tabular}{lcc}
\toprule
\textbf{Metric} & \textbf{Original} & \textbf{Box-Cox} \\
\midrule
Total features tested & 5 & 5 \\
Normality assumption passed & 0 & 2 \\
Homogeneity assumption passed & 3 & 4 \\
\textbf{All assumptions passed} & \textbf{0} & \textbf{2} \\
\bottomrule
\end{tabular}
\end{table}

\section{Detailed Comparison by Feature}

\begin{table}[h]
\centering
\caption{Feature-by-Feature Comparison: Original vs Box-Cox Transformed}
\begin{tabular}{p{4cm}ccc}
\toprule
\textbf{Feature} & \textbf{Original} & \textbf{Box-Cox} & \textbf{Result} \\
\midrule
Temperature (\textdegree{}C) & \textcolor{red}{FAIL} & \textcolor{red}{FAIL} & Same \\
Depth (m) & \textcolor{red}{FAIL} & \textcolor{green}{PASS} & \textcolor{blue}{Improved} \\
DO Bottom (mg/L) & \textcolor{red}{FAIL} & \textcolor{red}{FAIL} & Same \\
LOI (\%) & \textcolor{red}{FAIL} & \textcolor{green}{PASS} & \textcolor{blue}{Improved} \\
MPS (Phi) & \textcolor{red}{FAIL} & \textcolor{red}{FAIL} & Same \\

\bottomrule
\end{tabular}
\end{table}

\section{Box-Cox Transformation Details}

The Box-Cox transformation is defined as:
\[
y(\lambda) = 
\begin{cases}
\dfrac{x^\lambda - 1}{\lambda} & \text{if } \lambda \neq 0 \\[6pt]
\log(x) & \text{if } \lambda = 0
\end{cases}
\]

\subsection{Optimal Lambda Values and Interpretations}

\begin{itemize}
    \item \textbf{Temperature (\textdegree{}C)}: $\lambda$ = 1.9019 --- Custom power transformation
    \item \textbf{Measured Depth (m)}: $\lambda$ = 0.2392 --- Custom power transformation
    \item \textbf{Water DO Bottom (mg/L)}: $\lambda$ = 2.6106 --- Custom power transformation
    \item \textbf{LOI (\%)}: $\lambda$ = 0.0596 --- $\approx$ Log transformation
    \item \textbf{MPS (Phi)}: $\lambda$ = 3.3036 (shifted by 2.44) --- Custom power transformation

\end{itemize}

\section{Methodology Notes}

\begin{itemize}
    \item \textbf{Normality}: Assessed using Shapiro-Wilk test on ANOVA residuals (p $>$ 0.05 indicates normality)
    \item \textbf{Homogeneity of Variance}: Assessed using Levene's test (p $>$ 0.05 indicates homogeneous variances)
    \item \textbf{Independence}: Assumed based on spatial separation of sampling sites
    \item \textbf{Box-Cox Method}: Automatically finds optimal $\lambda$ to maximize normality of the transformed data
    \item \textbf{Shifting}: Features with non-positive values were shifted by $|\text{min}| + 1$ before transformation
    \item \textbf{Groups}: 4 taxa assemblage groups identified by Ward's hierarchical clustering
\end{itemize}

\end{document}
